% Statement and cited-input packet. Target proof is separately embargoed.
% Frozen II source SHA256: 74497bfe327fae25c9cdcb17334b3bf42a9793a9df7c0a4404d97568f2f760f0
% Target owner: cdp:st:B2; tag: 01BL

\begin{proposition}[A projective example with twelve length-one contributions]\hypertarget{tag-01BL}{}\label{cdp:st:B2}


Let $h:E\to\PP^1_t$ be the Weierstrass surface defined on the finite chart by
\[
 y^2z=x^3+(t^4+1)xz^2+t^6z^3,                          \tag{B.3}\hypertarget{tag-008I}{}\label{cdp:eq:B.3}
\]
and at $s=1/t$ by $x_\infty=s^2x$, $y_\infty=s^3y$, $z_\infty=z$. Let $O=[0:1:0]$ and $P=(0,t^3,1)$ be its two sections and put
\[
 M_E=\OO_E(P-O),\quad Y=E\times\PP^1,\quad
 g=h\circ\operatorname{pr}_E,\quad
 L=K_Y\otimes\operatorname{pr}_E^*M_E^{-1},\quad
 \mathcal E=T_Y\otimes L.
\]
The surface $E$ is smooth and projective. Its singular fibres are twelve irreducible nodal cubics, at the distinct roots $t_i$ of
\[
 D(t)=31t^{12}+12t^8+12t^4+4.
\]
For every such fibre $S_i=g^{-1}(t_i)$, with $A=(K_Y\otimes L^{-1})|_{S_i}$, we have
\[
H^0(S_i,N^*_{S_i/Y}\otimes A)=0,
\qquad H^0(S_i,\widetilde\Omega^1_{S_i}\otimes A)=0,
\]
\[
 (h^0,h^1,h^2)(S_i,\mathcal E|_{S_i})=(0,1,1).
\]
As coherent sheaves on the base, with noncanonical identifications of their one-dimensional fibres,
\[
 R^qg_*\mathcal E=0\ (q\ne2),\qquad
 R^2g_*\mathcal E=\bigoplus_{i=1}^{12}\C_{t_i}.        \tag{B.4}\hypertarget{tag-008J}{}\label{cdp:eq:B.4}
\]
In particular $(h^0,h^1,h^2,h^3)(Y,\mathcal E)=(0,0,12,0)$.
\end{proposition}

% Cited input cdp:st:B1; statement or labelled formula only.
\begin{proposition}[A complete nodal-family computation]\hypertarget{tag-01BK}{}\label{cdp:st:B1}

\label{wbII:thm:nodal-test}\label{wbII:prop:family-geometry}\label{wbII:prop:nodal-counterexample}
Let $0<\varepsilon<4/27$, let $\Delta=\{|t|<\varepsilon\}$, and define
\[
 \mathcal C=\{y^2z=x^2(x+z)+tz^3\}\subset\PP^2\times\Delta,
 \qquad X=\mathcal C\times\PP^1.
\]
The projection $X\to\Delta$ is proper, its total space is smooth, and its central fibre is $S=C_0\times\PP^1$, where $C_0$ is an irreducible nodal cubic. For each $\lambda\in\C^*$ let $A_\lambda$ be the line bundle whose pullback to the normalization $\nu:\PP^1_u\to C_0$ is trivial and whose sections in this trivialization satisfy $h(1)=\lambda h(-1)$. Put $A=\operatorname{pr}_1^*A_\lambda$ and $d_\lambda=1$ for $\lambda=1$, and $d_\lambda=0$ otherwise. In the order $q=0,1,2$, the coherent cohomology dimensions are
\[
\begin{array}{c|c}
\text{sheaf on }S & (h^0,h^1,h^2)\\\hline
 N^*_{S/X}\otimes A &(d_\lambda,d_\lambda,0)\\
 K_S\otimes A &(1,0,0)\\
 T_S\otimes A &(1,0,0)\\
 \widetilde\Omega^1_S\otimes A &(0,1+d_\lambda,d_\lambda)\\
 \Omega^1_S\otimes A &(1,1+d_\lambda,d_\lambda)\\
 \Omega^1_X|_S\otimes A &(1,1+d_\lambda,d_\lambda).
\end{array}                                                    \tag{B.1}\hypertarget{tag-008G}{}\label{cdp:eq:B.1}
\]
Here $T_S$ denotes the torsion of the cotangent sheaf, and $K_S$ its inverse image in the ambient cotangent sheaf. In particular the two proposed vanishings hold exactly for $\lambda\ne1$, and the ambient space of sections has dimension one for every $\lambda\in\C^*$. After shrinking $\Delta$, every $A_\lambda$ extends to a line bundle on $\mathcal C$ of the form $\OO_{\mathcal C}(P-Q)$ for holomorphic sections $P,Q$.
\end{proposition}

